\documentclass[11pt,a4paper]{article}
\usepackage[utf8]{inputenc}
\usepackage[T1]{fontenc}
\usepackage{amsmath,amsfonts,amssymb}
\usepackage{graphicx}
\usepackage{hyperref}
\usepackage{listings}
\usepackage{color}
\usepackage{booktabs}
\usepackage{float}
\usepackage{geometry}
\geometry{margin=1in}
\definecolor{zenblue}{RGB}{41,121,255}
\hypersetup{colorlinks=true,linkcolor=zenblue,urlcolor=zenblue,citecolor=zenblue}

\title{\textbf{Zen3-Embedding: A Packaging of Qwen3-Embedding\\
for Zen Retrieval}\\
\large Technical Whitepaper v2025.06}
\author{Zach Kelling \\ Zen LM Research Team\\
\texttt{research@zenlm.org}\\
\href{https://papers.zenlm.org}{papers.zenlm.org}}
\date{June 2025}

\begin{document}
\maketitle

\begin{abstract}
Zen3-Embedding is a repackaging of Alibaba's \textbf{Qwen3-Embedding} series, used as the
first-stage retrieval encoder in Zen retrieval pipelines. It is \emph{not} a from-scratch
model: it is the openly released, Apache-2.0 licensed \texttt{Qwen/Qwen3-Embedding} family
(0.6B, 4B, and 8B), produced by the Qwen team and built on the Qwen3 dense foundation
models~\cite{qwen3embedding}. Qwen3-Embedding is a \emph{causal} (decoder-only) LLM adapted
for embeddings: an \texttt{[EOS]} token is appended to the input and the final-layer hidden
state at that position is taken as the embedding. The models accept task instructions on the
query side, support Matryoshka-style user-defined output dimensions (from 32 up to the model's
native dimension), span 100+ languages including code, and handle a 32K context. Native
embedding dimensions are 1024 (0.6B), 2560 (4B), and 4096 (8B). On the Qwen team's evaluation,
the flagship Qwen3-Embedding-8B reports a Mean(Task) of \textbf{70.58} on the MTEB
Multilingual (MMTEB) leaderboard --- ranked No.~1 as of 5~June~2025 --- with MTEB English~v2
Mean(Task) 75.22, C-MTEB Mean(Task) 73.83, and MTEB-Code 80.68~\cite{qwen3embedding}. This
document describes how the upstream model is wired into Zen retrieval; all weights, training,
and reported numbers are the Qwen team's and are cited as such. Claims in earlier revisions of
a bespoke ``7680-dimensional'' bidirectional ``Zen MoDE'' encoder trained for ``Decentralized
Semantic Optimization,'' with BitDelta index compression, were fabricated and have been
removed.
\end{abstract}

\tableofcontents
\newpage

%% -----------------------------------------------------------------------
\section{Introduction}
\label{sec:intro}

Dense vector embeddings are the foundation of modern semantic search, retrieval-augmented
generation (RAG), and knowledge-base indexing. The quality of embeddings determines the
precision of document retrieval, which directly bounds the quality of RAG-powered
applications. Yet most deployed embedding models face a trilemma:

\begin{itemize}
\item \textbf{Quality vs. size}: High-quality embeddings require large models, but API
      latency and memory constraints favor small models.
\item \textbf{Generality vs. specialization}: General-purpose embeddings perform
      adequately on all tasks but rarely excel on any specific domain.
\item \textbf{Monolingual vs. multilingual}: Multilingual models traditionally
      underperform language-specific models on high-resource languages.
\end{itemize}

Rather than train an embedding model from scratch, Zen3-Embedding adopts the Qwen3-Embedding
series~\cite{qwen3embedding}, which targets these trade-offs directly. Our contribution is
integration and packaging, not modeling. The upstream model provides:

\begin{enumerate}
\item \textbf{A spectrum of sizes} --- 0.6B, 4B, and 8B, letting a deployment trade quality
      against latency and memory.
\item \textbf{Instruction-aware embeddings} --- The query may be prefixed with a task
      instruction that steers the representation toward the task, without fine-tuning.
\item \textbf{Matryoshka output dimensions} --- A single model supports user-defined output
      dimensions from 32 up to its native dimension, so index storage can be traded against
      retrieval quality without a separate model.
\item \textbf{Multilingual and code coverage} --- 100+ languages, including programming
      languages, with a 32K context.
\end{enumerate}

We do not modify the Qwen3-Embedding weights; Zen3-Embedding is the upstream model wired into
the Zen retrieval stack.

%% -----------------------------------------------------------------------
\section{Architecture}
\label{sec:arch}

\subsection{Causal Decoder Adapted for Embeddings}

Qwen3-Embedding is \emph{not} a bidirectional BERT-style encoder and has no
mixture-of-experts ``Zen MoDE'' backbone. Per the Qwen team~\cite{qwen3embedding}, each model
is a fine-tune of the corresponding Qwen3-\emph{Base} \emph{dense, causal} (decoder-only) LLM
(e.g. the 8B from \texttt{Qwen3-8B-Base}, the 0.6B from \texttt{Qwen3-0.6B-Base}). The
``7.1B-parameter, 48-layer, 64-expert top-4 bidirectional'' specification in earlier revisions
did not correspond to any deployed weights and has been removed.

Released sizes and native embedding dimensions~\cite{qwen3embedding}:
\begin{itemize}
\item Qwen3-Embedding-0.6B: 28 layers, native embedding dimension 1024.
\item Qwen3-Embedding-4B: 36 layers, native embedding dimension 2560.
\item Qwen3-Embedding-8B: 36 layers, native embedding dimension 4096.
\item Context length: 32,768 tokens (all sizes).
\end{itemize}

\subsection{Embedding Extraction (Last-Token Pooling)}

The embedding is the final-layer hidden state at an appended \texttt{[EOS]}
token~\cite{qwen3embedding}. Because attention is causal, that last position has attended over
the entire input, so its hidden state serves as the sequence representation:

\begin{equation}
\mathbf{e} = \mathbf{h}^{(L)}_{\texttt{[EOS]}} \in \mathbb{R}^{d_{\text{model}}},
\end{equation}

where $d_{\text{model}}$ is the model's native dimension (1024 / 2560 / 4096). There is no
\texttt{[CLS]} token, no mean$+$CLS dual pooling, and no learned projection to a 7680-dim
space; those constructions in earlier revisions were fabricated.

\subsection{Matryoshka Output Dimensions}

Qwen3-Embedding supports user-defined output dimensions from 32 up to the native
dimension~\cite{qwen3embedding,mrl}, so a shorter prefix of the embedding can be used to
shrink index storage at a modest quality cost. The supported range is continuous (32 to
$d_{\text{model}}$), not the fixed ladder $\{256,\dots,7680\}$ asserted in earlier revisions.

\subsection{Use in the Zen Retrieval Stack}

Zen indexing computes Qwen3-Embedding vectors for documents (no instruction) and builds a
standard vector index (e.g. HNSW) over them; at query time the query embedding (optionally
with a task instruction) is used for similarity search, and the Qwen3-Reranker stage
re-scores the shortlist. The ``DSO-native 7680-dim format,'' ``BitDelta 1.58-bit index
compression,'' and ``peer-to-peer node'' descriptions of earlier revisions did not describe
this model and have been removed.

\subsection{Instruction-Aware Embedding}

Task instructions are prepended to queries using a structured template:

\begin{verbatim}
Instruction: <task_description>
Query: <query_text>
\end{verbatim}

The instruction shifts the attention pattern toward task-relevant dimensions of the
embedding space. Evaluated instruction types include:
\begin{itemize}
\item \texttt{Retrieve relevant documents}: general retrieval
\item \texttt{Find semantically similar sentences}: semantic similarity
\item \texttt{Classify the following text}: classification
\item \texttt{Find duplicates}: deduplication
\end{itemize}

Instructions are applied to queries only (not documents), preserving document index reusability.

%% -----------------------------------------------------------------------
\section{Training}
\label{sec:training}

\subsection{Provenance, Not Training}

Zen3-Embedding performs no training of its own. The models were trained by the Qwen team; we
summarize their published approach at a high level and attribute it
accordingly~\cite{qwen3embedding}.

Per the Qwen team, Qwen3-Embedding starts from the Qwen3-Base dense LLMs and is adapted to
embeddings with a contrastive objective (the standard last-token / \texttt{[EOS]}
representation with an InfoNCE-style loss over positive (query, document) pairs and
negatives), with the Qwen3 generative models used to synthesize a large, multilingual portion
of the training pairs~\cite{qwen3embedding}. The contrastive loss has the usual form

\begin{equation}
\mathcal{L}_{\text{InfoNCE}} = -\log \frac{\exp(\text{sim}(\mathbf{q}, \mathbf{d}^+) / \tau)}{\exp(\text{sim}(\mathbf{q}, \mathbf{d}^+) / \tau) + \sum_{j} \exp(\text{sim}(\mathbf{q}, \mathbf{d}^-_j) / \tau)},
\end{equation}

with $\text{sim}(\cdot,\cdot)$ cosine similarity. For the authoritative dataset composition,
synthetic-data pipeline, and hyperparameters, see the Qwen team's report~\cite{qwen3embedding}.

The fabricated training details in earlier revisions --- ``bidirectional MLM adaptation,'' a
2.1B-pair corpus with specific per-source counts, a synthetic generator named
``Zen3-Medium,'' and a 32{,}768 GradCache batch --- did not describe any real training run
and have been removed.

%% -----------------------------------------------------------------------
\section{Evaluation}
\label{sec:eval}

All numbers in this section are the Qwen team's published results for
Qwen3-Embedding~\cite{qwen3embedding}, reproduced here with attribution. We ran no evaluations
of our own and make no independent benchmark claims.

\subsection{Flagship (8B) Headline Numbers}

\begin{table}[H]
\centering
\caption{Qwen3-Embedding-8B, as reported by the Qwen team~\cite{qwen3embedding}. The MMTEB
Mean(Task) of 70.58 ranked No.~1 on the MTEB Multilingual leaderboard as of 5~June~2025.}
\begin{tabular}{lc}
\toprule
\textbf{Benchmark} & \textbf{Score} \\
\midrule
MTEB Multilingual (MMTEB), Mean(Task) & 70.58 \\
MTEB Multilingual (MMTEB), Mean(Type) & 61.69 \\
MTEB English v2, Mean(Task)           & 75.22 \\
MTEB English v2, Mean(Type)           & 68.70 \\
C-MTEB (Chinese), Mean(Task)          & 73.83 \\
MTEB-Code (nDCG@10)                   & 80.68 \\
\bottomrule
\end{tabular}
\label{tab:qwen3emb-8b}
\end{table}

\subsection{Across Model Sizes}

\begin{table}[H]
\centering
\caption{Released sizes and native embedding dimensions~\cite{qwen3embedding}. The 8B is the
flagship; the Qwen team's report gives the full MMTEB breakdown for each size, with quality
increasing with size.}
\begin{tabular}{lcc}
\toprule
\textbf{Model} & \textbf{Layers} & \textbf{Native dim} \\
\midrule
Qwen3-Embedding-0.6B & 28 & 1024 \\
Qwen3-Embedding-4B   & 36 & 2560 \\
Qwen3-Embedding-8B   & 36 & 4096 \\
\bottomrule
\end{tabular}
\label{tab:qwen3emb-sizes}
\end{table}

The 8B MMTEB Mean(Task) of 70.58 is given above; per-size MMTEB/MTEB scores for the 0.6B and
4B models are reported in the Qwen team's paper and we refer the reader there rather than
restating numbers we have not independently verified~\cite{qwen3embedding}.

\subsection{A Note on Earlier Revisions}

Earlier revisions of this document reported a fabricated ``MTEB average 74.3,'' a ``Zen2-Embedding''
comparison column, BEIR/MIRACL tables, and a Matryoshka/BitDelta storage table at 7680
dimensions. None of those described this model or any real measurement; they were fabricated
and have been removed. For per-task and per-dataset breakdowns of the genuine numbers, see the
Qwen team's report and the HuggingFace model cards~\cite{qwen3embedding}.

%% -----------------------------------------------------------------------
\section{Related Work}
\label{sec:related}

\subsection{Dense Retrieval}

Dense passage retrieval~\cite{karpukhin2020dense} established bi-encoder architectures
as practical alternatives to BM25 for open-domain question answering. Qwen3-Embedding sits in
the more recent line of repurposing decoder-only LLMs as embedding models via last-token
pooling and contrastive fine-tuning~\cite{qwen3embedding}.

\subsection{Contrastive Learning for Embeddings}

SimCSE~\cite{gao2021simcse} demonstrated that contrastive learning with dropout as
data augmentation substantially improves semantic textual similarity. E5~\cite{wang2022e5}
and subsequent work showed that web-scale weakly supervised contrastive training
with curated hard negatives matches or exceeds expensive human-annotated pairs.
Qwen3-Embedding follows this contrastive-fine-tuning paradigm on top of the Qwen3 base
LLMs~\cite{qwen3embedding}.

\subsection{Matryoshka Representations}

Matryoshka Representation Learning~\cite{mrl} enables flexible dimension truncation
with minimal quality loss; Qwen3-Embedding exposes this as user-selectable output dimensions
from 32 up to each model's native dimension~\cite{qwen3embedding}.

%% -----------------------------------------------------------------------
\section{Conclusion}
\label{sec:conclusion}

Zen3-Embedding is the first-stage retrieval encoder of Zen retrieval pipelines, implemented as
a direct deployment of Alibaba's Apache-2.0 Qwen3-Embedding models
(0.6B/4B/8B)~\cite{qwen3embedding}. Qwen3-Embedding is a causal decoder LLM adapted for
embeddings via last-token (\texttt{[EOS]}) pooling, with instruction-aware queries, Matryoshka
output dimensions, 100+ language coverage, and a 32K context. The flagship 8B reports an MMTEB
Mean(Task) of 70.58 (No.~1 as of 5~June~2025) per the Qwen team. We contribute integration into
the Zen retrieval stack, not a new model, and we report the Qwen team's published numbers with
attribution rather than benchmarks of our own. The ``DSO-native 7680-dimensional Zen MoDE
encoder'' and ``BitDelta index compression'' claims of earlier revisions have been removed as
fabrications.

\section*{Acknowledgments and Attribution}

Zen3-Embedding is built entirely on Qwen3-Embedding (\texttt{Qwen/Qwen3-Embedding},
Apache-2.0), developed by the Qwen team at Alibaba; all model weights, training, and reported
metrics are theirs. We thank the Qwen team for releasing the models openly and the MTEB/MMTEB
maintainers for the evaluation suites.

\bibliographystyle{plain}
\begin{thebibliography}{9}

\bibitem{karpukhin2020dense}
V. Karpukhin et al.,
``Dense Passage Retrieval for Open-Domain Question Answering,''
\textit{EMNLP}, 2020.

\bibitem{gao2021simcse}
T. Gao, X. Yao, D. Chen,
``SimCSE: Simple Contrastive Learning of Sentence Embeddings,''
\textit{EMNLP}, 2021.

\bibitem{wang2022e5}
L. Wang et al.,
``Text Embeddings by Weakly-Supervised Contrastive Pre-Training,''
\textit{arXiv:2212.03533}, 2022.

\bibitem{mrl}
A. Kusupati et al.,
``Matryoshka Representation Learning,''
\textit{NeurIPS}, 2022.

\bibitem{qwen3embedding}
Y. Zhang, M. Li, D. Long, X. Zhang, H. Lin, B. Yang, P. Xie, A. Yang, D. Liu, J. Lin,
F. Huang, J. Zhou (Qwen Team, Alibaba),
``Qwen3 Embedding: Advancing Text Embedding and Reranking Through Foundation Models,''
\textit{arXiv:2506.05176}, 2025.
Models: \texttt{Qwen/Qwen3-Embedding-\{0.6B,4B,8B\}} (Apache-2.0).

\end{thebibliography}

\end{document}
